%% ============================================================================
%%  Trust the Figure, Not the Caption
%%  Conference paper scaffold -- IEEEtran conference format
%%
%%  STATUS: DRAFT SCAFFOLD.
%%  Sections I-IV and VI contain real prose grounded in the implemented system.
%%  Section V (Results) is intentionally EMPTY -- fill only from executed runs.
%%  Every \TODO marker below is a blocker for submission.
%%
%%  Compile:  pdflatex main && bibtex main && pdflatex main && pdflatex main
%% ============================================================================

\documentclass[conference]{IEEEtran}
\IEEEoverridecommandlockouts

\usepackage{cite}
\usepackage{amsmath,amssymb}
\usepackage{algorithmic}
\usepackage{graphicx}
\usepackage{textcomp}
\usepackage{xcolor}
\usepackage{booktabs}
\usepackage{multirow}
\usepackage{url}

% ---- Draft-mode TODO macro. Set \draftmodefalse before submitting; any
% ---- remaining TODO will then vanish silently, so grep for \TODO first.
\newif\ifdraftmode
\draftmodetrue
\newcommand{\TODO}[1]{\ifdraftmode\textcolor{red}{\textbf{[TODO: #1]}}\fi}
\newcommand{\NUM}[1]{\ifdraftmode\textcolor{red}{\textbf{[#1]}}\else#1\fi}

\begin{document}

\title{Trust the Figure, Not the Caption: Visual--Textual Consistency
Checking in Multimodal Retrieval-Augmented Generation over Research Papers}

\author{
\IEEEauthorblockN{\TODO{author list}}
\IEEEauthorblockA{\textit{Department of Computer Science and Engineering} \\
\TODO{institution} \\
\TODO{emails}}
}

\maketitle

% ============================================================================
\begin{abstract}
Retrieval-augmented generation (RAG) over scientific literature is typically
evaluated on its ability to locate and reproduce textual evidence. Yet a
substantial fraction of the claims in a research paper are carried by figures
and tables, and the caption that describes a figure is not always faithful to
what the figure shows. We present a hybrid multimodal RAG system for
question answering over a single research paper, in which figures and tables
are preserved as first-class retrievable evidence and passed to a multimodal
language model alongside the text. Building on this system, we introduce a
\emph{caption-perturbation protocol} that measures whether a multimodal RAG
system notices when a caption contradicts the figure it describes, and whether
it resolves the conflict in favour of the visual evidence. We evaluate
\NUM{$N$} configurations over a corpus of \NUM{$M$} annotated research papers.
\TODO{one-sentence result statement; write last, from real numbers}
\end{abstract}

\begin{IEEEkeywords}
retrieval-augmented generation, multimodal retrieval, hybrid search,
scientific document understanding, visual question answering, groundedness
\end{IEEEkeywords}

% ============================================================================
\section{Introduction}

Question answering over research papers is a natural application of
retrieval-augmented generation~\cite{lewis2020rag}: a reader arrives with a
specific question, the relevant evidence is confined to a single document, and
the cost of a fabricated answer is high. Most deployed systems address this by
chunking the document's text, retrieving the top-scoring chunks, and
conditioning a language model on them.

This text-centric view discards a large part of what a research paper actually
communicates. Results are reported in tables; trends, architectures, and
ablations are reported in figures. A pipeline that flattens a PDF into a
sequence of text spans either loses this content entirely or reduces it to the
caption --- a one- or two-sentence summary written by the authors.

Relying on the caption introduces a failure mode that, to our knowledge, has
not been measured directly in the RAG literature. A caption is a
\emph{claim about} a figure, and claims can be wrong. Captions are copy-edited
late, figures are regenerated after captions are written, and sign errors,
reversed trends, and stale magnitudes survive review. A RAG system that reads
only the caption inherits every such error with no means of detecting it. A
system that reads the figure \emph{and} the caption has, in principle, the
evidence needed to notice the discrepancy --- but only if it is induced to
check.

We investigate this directly. Our contributions are:

\begin{enumerate}
  \item \textbf{A hybrid multimodal RAG pipeline} for single-paper QA
  (Section~\ref{sec:system}) that parses figures and tables into cropped
  image evidence with structured provenance, retrieves over text and captions
  using reciprocal rank fusion of lexical and dense rankings, and links
  retrieved text back to the figures it references before generation.

  \item \textbf{A caption-perturbation protocol} (Section~\ref{sec:protocol})
  for measuring visual--textual consistency checking. We generate controlled
  contradictions between a figure and its caption and measure both the
  detection rate under perturbation and, critically, the false-alarm rate on
  unperturbed controls --- a quantity that a detection-only evaluation would
  hide.

  \item \textbf{An empirical study} over \NUM{$M$} annotated papers
  isolating the contribution of visual evidence and of explicit
  consistency-checking instructions.
  \TODO{finalise claim once Section V exists}
\end{enumerate}

\TODO{Add a Figure 1: side-by-side of a figure whose caption contradicts it,
with the system's flagged output. This is the paper's hook and reviewers will
look for it on page 1.}

% ============================================================================
\section{Related Work}

\subsection{Hybrid Retrieval}
Lexical retrieval with BM25~\cite{robertson2009bm25} and dense retrieval with
learned bi-encoders~\cite{karpukhin2020dpr} have complementary failure modes:
the former is precise on rare terms and brittle to paraphrase, the latter
robust to paraphrase and prone to missing exact identifiers. Reciprocal rank
fusion (RRF)~\cite{cormack2009rrf} combines ranked lists without score
normalisation or tuned weights, and remains a strong default. We adopt RRF as
infrastructure rather than as a contribution.

\subsection{Multimodal Document Retrieval}
Recent work retrieves over rendered document pages directly, bypassing text
extraction: ColPali~\cite{faysse2024colpali} embeds page images with a
vision--language model and matches them against queries at the patch level.
Such approaches sidestep parsing error but forgo the structured provenance
(section, page, element type) that a parser recovers. Our pipeline takes the
parsing route, using Docling~\cite{docling2024} for layout analysis, and
retains figures as separately addressable evidence.
\TODO{Position against page-level retrieval explicitly --- a reviewer will ask
why you parse instead of embedding pages. The honest answer is provenance and
per-element citation; make sure that is argued, not asserted.}

\subsection{Groundedness and Hallucination}
Evaluation frameworks such as RAGAS~\cite{es2023ragas} score whether a
generated answer is supported by retrieved context, and hallucination
surveys~\cite{ji2023hallucination} catalogue the ways generation departs from
evidence. These frameworks assume the retrieved context is itself correct. The
setting we study is different: the retrieved context is
\emph{internally inconsistent}, and the system must adjudicate between two
conflicting pieces of evidence rather than remain faithful to a single
consistent one.
\TODO{Search for prior work on chart--caption factual consistency and on
scientific claim verification against figures; if a close neighbour exists,
this subsection must engage it directly rather than claim novelty.}

% ============================================================================
\section{System Architecture}
\label{sec:system}

The system processes one paper at a time in four stages: parsing, indexing,
evidence linking, and multimodal generation.

\subsection{Parsing and Chunking}
A PDF is converted with Docling~\cite{docling2024} configured to emit images
for both picture and table elements. We iterate the resulting document tree and
emit one chunk per element, typed as \texttt{text}, \texttt{table}, or
\texttt{figure}. Each chunk carries: a unique identifier; the source page
number; the enclosing section heading, tracked by carrying forward the most
recent section header encountered during traversal; the element text; and, for
non-text elements, the caption, the concatenation of the three preceding text
elements as \texttt{nearby\_text}, and a path to a cropped image.

Images are obtained from Docling's native element rasteriser where available.
Where it is not, we fall back to cropping the element's bounding box from the
source PDF with PyMuPDF at 150 DPI. Table elements additionally retain a
Markdown serialisation of their cell structure, so that a table is represented
both as searchable text and as an image.

\subsection{Hybrid Retrieval}
Both retrievers index the same derived string per chunk: the section heading,
the element text, and the caption, concatenated. This makes figures and tables
retrievable through their captions and surrounding section context even though
the retrievers themselves operate on text.

\textbf{Lexical.} A BM25 index~\cite{robertson2009bm25} over whitespace-split,
lowercased tokens.

\textbf{Dense.} Chunk strings are embedded with
\texttt{BAAI/bge-small-en-v1.5}~\cite{xiao2023cpack} and stored in a persistent
ChromaDB collection under cosine distance; retrieval scores are reported as
$1 - d_{\cos}$.

\textbf{Fusion.} We retrieve the top 60 candidates from each retriever and
combine them by reciprocal rank fusion~\cite{cormack2009rrf}. For a chunk $c$
appearing at rank $r_i(c)$ in ranked list $i$,
\begin{equation}
\mathrm{RRF}(c) = \sum_{i} \frac{1}{k + r_i(c)}, \qquad k = 60,
\end{equation}
with chunks absent from a list contributing nothing. The fused list is sorted
by $\mathrm{RRF}(c)$ and truncated to the top 5 for generation.

\subsection{Evidence Linking}
Retrieval returns text chunks, but a question about a trend is answered by the
figure that the text refers to. After fusion we scan each retrieved text chunk
for references of the form \emph{Figure~$n$} / \emph{Fig.~$n$} /
\emph{Table~$n$}, and attach any figure or table chunk whose caption matches
the referenced ordinal and whose element type agrees. Attached images are
carried into the generation context alongside the text chunk that referenced
them.

This mechanism is deliberately simple, and its coverage is bounded by caption
availability. In a single-paper pilot, parsing yielded 117 chunks --- 97 text,
16 figure, 4 table --- of which 20 carried a cropped image but only 10 carried
a caption. Half of the extracted visual elements therefore cannot be reached by
caption-based linking at all.
\TODO{Recompute this statistic across the full annotated corpus. Reported here
from one paper; label it as a pilot number or replace it. Do not let a
single-document count appear in the final paper without that qualification.}

\subsection{Multimodal Generation}
The top-ranked chunks and their linked images are assembled into a single
multimodal prompt. Text chunks are serialised with their index, page number,
and section, so that the model can cite evidence positionally; images are
attached as raw image parts. Generation uses a multimodal LLM at temperature
0.2.
\TODO{Fix the exact model identifier and version string, and verify it against
the provider's model list before submission. The identifier currently in the
codebase should not be copied into the paper unchecked.}

The system instruction directs the model to answer from the provided evidence
and to cite context indices. It additionally carries the
\emph{consistency-checking directive} that is the object of study in
Section~\ref{sec:protocol}: the model is told to compare each figure against
its caption and surrounding text, to flag any contradiction explicitly, and to
resolve conflicts in favour of the visual evidence. Section~\ref{sec:protocol}
treats this directive as an ablatable condition rather than a fixed property of
the system.

% ============================================================================
\section{Experimental Design}
\label{sec:protocol}

\subsection{Corpus and Annotation}
We assemble a corpus of \NUM{$M$} research papers
\TODO{state the source venue/arXiv categories and the sampling procedure; a
convenience sample must be described as one}.
For each paper, annotators author question--answer pairs and mark the set of
gold chunks that support each answer.

\textbf{Identifier stability.} Gold annotations reference chunks by
identifier, so identifiers must be reproducible across re-parses. We assign
each chunk a content-derived identifier --- a hash over
(document, page, element type, element text) --- rather than a random one, so
that an annotation set remains valid when the corpus is re-parsed.
\TODO{This is a requirement on the implementation, not a description of it.
The current parser assigns \texttt{uuid4()} per chunk, which invalidates gold
labels on every re-parse. Change the parser before annotating anything, or the
annotation effort is wasted.}

\subsection{Retrieval Evaluation}
We report Recall@$k$, MRR, and nDCG@$k$ over the annotated gold sets for three
configurations: BM25 alone, dense retrieval alone, and RRF fusion. This
comparison characterises the retrieval substrate; it is not the paper's claim.

\TODO{The evaluation harness currently substitutes the fusion system's own
top-ranked output for the gold labels when the annotated identifiers do not
resolve. Under that substitution the fusion condition scores perfectly by
construction and the baselines are scored against fusion's output rather than
against ground truth. No number produced by the harness in its present state
may enter this paper. Remove the substitution and make unresolvable gold
identifiers a hard error before running anything.}

\subsection{Caption-Perturbation Protocol}
The central experiment measures whether the system detects figure--caption
contradictions.

\textbf{Perturbations.} For each figure with a caption, we construct a
contradicted variant by a single controlled edit of one of four types:
\begin{itemize}
  \item \emph{Trend reversal}: increasing $\leftrightarrow$ decreasing,
  improves $\leftrightarrow$ degrades.
  \item \emph{Relational inversion}: outperforms $\leftrightarrow$
  underperforms, higher $\leftrightarrow$ lower.
  \item \emph{Magnitude alteration}: a reported value replaced with one
  inconsistent with the plotted data.
  \item \emph{Entity substitution}: an axis label, method name, or series
  identifier swapped for another present in the figure.
\end{itemize}
Each perturbed caption remains fluent and plausible in isolation; the
contradiction is detectable only by consulting the figure.
\TODO{Perturbations must be human-verified as (a) genuinely contradictory and
(b) undetectable from text alone. Report inter-annotator agreement on a
sample. A reviewer will assume the perturbations leak textual cues unless you
show otherwise.}

\textbf{Conditions.} We cross perturbation (perturbed vs.\ unperturbed
control) with three system configurations:
\begin{enumerate}
  \item \textbf{Text-only}: captions and text in context, no images.
  \item \textbf{Multimodal, no directive}: images in context, consistency
  directive removed from the system instruction.
  \item \textbf{Multimodal + directive}: the full system.
\end{enumerate}
Condition~1 establishes that the contradiction is not textually detectable.
The 1$\rightarrow$2 contrast isolates the contribution of visual evidence; the
2$\rightarrow$3 contrast isolates the contribution of the instruction.

\textbf{Metrics.}
\begin{itemize}
  \item \emph{Detection rate}: proportion of perturbed items for which the
  system explicitly flags a contradiction.
  \item \emph{False-alarm rate}: proportion of \emph{unperturbed} controls
  falsely flagged. A system that flags indiscriminately achieves perfect
  detection; without this number the detection rate is uninterpretable.
  \item \emph{Resolution accuracy}: given a question whose answer differs
  under figure and caption, whether the answer follows the figure.
\end{itemize}
\TODO{Specify how flagging is adjudicated --- keyword match on the output is
too brittle to defend. Either define a rubric and have humans score a sample,
or use an LLM judge and report its agreement with human scores on that sample.
Do not report judge scores without that agreement number.}

% ============================================================================
\section{Results}
\label{sec:results}

\TODO{THIS SECTION IS EMPTY BY DESIGN. Populate only from executed runs.
Prerequisites, in order:
(1) content-derived chunk identifiers in the parser;
(2) ground-truth substitution removed from the evaluation harness;
(3) the annotated corpus built;
(4) the perturbation set constructed and human-verified;
(5) runs executed and logged.
Table skeletons below are placeholders --- they contain no data.}

\begin{table}[t]
\caption{Retrieval performance over the annotated corpus.}
\label{tab:retrieval}
\centering
\begin{tabular}{lccc}
\toprule
Configuration & Recall@5 & MRR & nDCG@5 \\
\midrule
BM25 only      & \TODO{} & \TODO{} & \TODO{} \\
Dense only     & \TODO{} & \TODO{} & \TODO{} \\
RRF fusion     & \TODO{} & \TODO{} & \TODO{} \\
\bottomrule
\end{tabular}
\end{table}

\begin{table}[t]
\caption{Caption-perturbation results. Detection is measured on perturbed
items, false alarms on unperturbed controls.}
\label{tab:perturbation}
\centering
\begin{tabular}{lccc}
\toprule
Configuration & Detection & False alarm & Resolution \\
\midrule
Text-only              & \TODO{} & \TODO{} & \TODO{} \\
Multimodal, no directive & \TODO{} & \TODO{} & \TODO{} \\
Multimodal + directive & \TODO{} & \TODO{} & \TODO{} \\
\bottomrule
\end{tabular}
\end{table}

\TODO{Report variance. Single-run numbers from a stochastic decoder are not
evidence; run each condition multiple times with fixed seeds where available
and report mean and spread, or state explicitly that decoding is deterministic
at this temperature and show it.}

% ============================================================================
\section{Discussion and Limitations}

\TODO{Write after Results. The items below are known in advance and must
appear regardless of what the numbers show.}

\textbf{Retrieval is text-mediated.} Figures are retrievable only through
their captions, section headings, and surrounding text; the pipeline never
embeds image content for retrieval. An uncaptioned figure is effectively
invisible to the retriever, and in our pilot parse half the extracted visual
elements were uncaptioned. This bounds every result we report.

\textbf{Evidence linking is regex-based.} Cross-references are matched by
ordinal against caption strings. References of the form ``the figure above'',
subfigure references (``Fig.~3(b)''), and figures discussed without explicit
numbering are not linked.

\textbf{Retrieval component configuration is untuned.} The lexical index uses
whitespace tokenisation with no stemming or stopword removal. The dense
retriever omits the asymmetric query instruction prefix that the BGE model
family is trained to expect, which is known to affect retrieval quality on
short queries. Both are defensible as defaults but neither is tuned, so the
retrieval comparison in Table~\ref{tab:retrieval} should be read as a
characterisation of this pipeline rather than as a general claim about lexical
versus dense retrieval.

\textbf{Chunk granularity.} One chunk is emitted per parsed document element,
without merging. Elements are consequently short, which favours lexical
matching on rare terms and may fragment arguments spanning several elements.

\textbf{Perturbations are synthetic.} Constructed contradictions are a proxy
for naturally occurring caption errors. They establish capability under a
controlled contradiction; they do not establish the rate at which such errors
occur in published work, nor that natural errors resemble synthetic ones.
\TODO{If feasible, include even a small set of naturally occurring
caption--figure discrepancies. A handful of real cases substantially
strengthens the paper's external validity.}

\textbf{Single-document scope.} The system indexes one paper per session.
Cross-paper retrieval is out of scope.

% ============================================================================
\section{Conclusion}

\TODO{Write last. State only what the experiments in Section V support. If the
directive ablation shows no effect, that is a publishable negative result about
prompt-level consistency checking --- report it as such rather than reframing
the paper around the retrieval comparison, which is not novel.}

% ============================================================================
\section*{Reproducibility}

\TODO{State the model identifier and version, embedding model revision,
retrieval hyperparameters ($k=60$ for both candidate depth and the RRF
constant), decoding temperature, and the date of API access. Release the
annotated corpus, the perturbation set, and the evaluation harness. Note that
results depend on a hosted model whose behaviour may change; the access date
is not optional.}

% ============================================================================
% NOTE: Every entry below must be verified against the actual publication
% before submission -- authors, venue, year, and page numbers. Do not submit
% a bibliography that has not been checked entry by entry.
% ============================================================================
\begin{thebibliography}{00}

\bibitem{lewis2020rag}
P.~Lewis \emph{et al.}, ``Retrieval-augmented generation for knowledge-intensive
NLP tasks,'' in \emph{Advances in Neural Information Processing Systems}, 2020.

\bibitem{robertson2009bm25}
S.~Robertson and H.~Zaragoza, ``The probabilistic relevance framework: BM25 and
beyond,'' \emph{Foundations and Trends in Information Retrieval}, vol.~3,
no.~4, pp.~333--389, 2009.

\bibitem{karpukhin2020dpr}
V.~Karpukhin \emph{et al.}, ``Dense passage retrieval for open-domain question
answering,'' in \emph{Proc. EMNLP}, 2020.

\bibitem{cormack2009rrf}
G.~V. Cormack, C.~L.~A. Clarke, and S.~Buettcher, ``Reciprocal rank fusion
outperforms Condorcet and individual rank learning methods,'' in
\emph{Proc. SIGIR}, 2009.

\bibitem{xiao2023cpack}
S.~Xiao \emph{et al.}, ``C-Pack: Packed resources for general Chinese
embeddings,'' arXiv preprint, 2023.

\bibitem{faysse2024colpali}
M.~Faysse \emph{et al.}, ``ColPali: Efficient document retrieval with vision
language models,'' arXiv preprint, 2024.

\bibitem{docling2024}
Docling Team, ``Docling technical report,'' IBM Research, arXiv preprint, 2024.

\bibitem{es2023ragas}
S.~Es \emph{et al.}, ``RAGAS: Automated evaluation of retrieval augmented
generation,'' arXiv preprint, 2023.

\bibitem{ji2023hallucination}
Z.~Ji \emph{et al.}, ``Survey of hallucination in natural language
generation,'' \emph{ACM Computing Surveys}, vol.~55, no.~12, pp.~1--38, 2023.

\end{thebibliography}

\end{document}
